\documentclass[12pt]{article}
\usepackage{latexblog}

% ============================================================
% Blog Metadata
% ============================================================
\blogtitle{阅读笔记：ConcurrentHashMap}
\blogdate{2020-03-19}
\blogtags{刨根问底, 并发编程, JDK1.8, Concurrent}
\bloglang{zh}
\blogsummary{}

\begin{document}
\makeblogtitle

我们知道HashMap不是Thread-safe的，而HashTable内部采取了同步操作，是线程安全的。然而有趣的是你去看HashTable的文档，它会建议你：如果不要Thread-Safe你就用HashMap吧，否则你用ConcurrentHashMap好了。

一般如果对线程安全有要求，我们有如下的一些选择：

\begin{itemize}
\tightlist
\item
  ConcurrentHashMap
\item
  Hashtable
\item
  Collections.synchronizedMap
\end{itemize}

\section{Collections.synchronizedMap}\label{collections.synchronizedmap}

这个实现很粗暴，实际上就是将Map的各个操作都进行了包装和同步：

\begin{Shaded}
\begin{Highlighting}[]
\KeywordTok{private} \DataTypeTok{static} \KeywordTok{class}\NormalTok{ SynchronizedMap}\OperatorTok{\textless{}}\NormalTok{K}\OperatorTok{,}\NormalTok{V}\OperatorTok{\textgreater{}}
        \KeywordTok{implements} \BuiltInTok{Map}\OperatorTok{\textless{}}\NormalTok{K}\OperatorTok{,}\NormalTok{V}\OperatorTok{\textgreater{},} \BuiltInTok{Serializable} \OperatorTok{\{}
    \KeywordTok{private} \DataTypeTok{final} \BuiltInTok{Map}\OperatorTok{\textless{}}\NormalTok{K}\OperatorTok{,}\NormalTok{V}\OperatorTok{\textgreater{}}\NormalTok{ m}\OperatorTok{;}     \CommentTok{// Backing Map}
    \DataTypeTok{final} \BuiltInTok{Object}\NormalTok{      mutex}\OperatorTok{;}  

    \FunctionTok{SynchronizedMap}\OperatorTok{(}\BuiltInTok{Map}\OperatorTok{\textless{}}\NormalTok{K}\OperatorTok{,}\NormalTok{V}\OperatorTok{\textgreater{}}\NormalTok{ m}\OperatorTok{)} \OperatorTok{\{}
        \KeywordTok{this}\OperatorTok{.}\FunctionTok{m} \OperatorTok{=}\NormalTok{ Objects}\OperatorTok{.}\FunctionTok{requireNonNull}\OperatorTok{(}\NormalTok{m}\OperatorTok{);}
\NormalTok{        mutex }\OperatorTok{=} \KeywordTok{this}\OperatorTok{;}
    \OperatorTok{\}}
\end{Highlighting}
\end{Shaded}

在构造函数中传入了原来的Map，以及一个对象锁（如果不传那就默认是this了）。然后，所有的操作都进行了同步处理：

\begin{Shaded}
\begin{Highlighting}[]
\KeywordTok{public} \DataTypeTok{boolean} \FunctionTok{containsValue}\OperatorTok{(}\BuiltInTok{Object}\NormalTok{ value}\OperatorTok{)} \OperatorTok{\{}
    \KeywordTok{synchronized} \OperatorTok{(}\NormalTok{mutex}\OperatorTok{)} \OperatorTok{\{}\ControlFlowTok{return}\NormalTok{ m}\OperatorTok{.}\FunctionTok{containsValue}\OperatorTok{(}\NormalTok{value}\OperatorTok{);\}}
\OperatorTok{\}}
\KeywordTok{public}\NormalTok{ V }\FunctionTok{get}\OperatorTok{(}\BuiltInTok{Object}\NormalTok{ key}\OperatorTok{)} \OperatorTok{\{}
    \KeywordTok{synchronized} \OperatorTok{(}\NormalTok{mutex}\OperatorTok{)} \OperatorTok{\{}\ControlFlowTok{return}\NormalTok{ m}\OperatorTok{.}\FunctionTok{get}\OperatorTok{(}\NormalTok{key}\OperatorTok{);\}}
\OperatorTok{\}}
\CommentTok{// ...}
\end{Highlighting}
\end{Shaded}

\section{HashTable}\label{hashtable}

HashTable实现线程安全的方式与上面有些类似，对所有需要同步的地方直接进行了同步：

\begin{Shaded}
\begin{Highlighting}[]
\KeywordTok{public} \KeywordTok{synchronized} \DataTypeTok{int} \FunctionTok{size}\OperatorTok{()} \OperatorTok{\{}
    \ControlFlowTok{return}\NormalTok{ count}\OperatorTok{;}
\OperatorTok{\}}
\end{Highlighting}
\end{Shaded}

那么HashMap和HashTable有什么区别呢？除了同步之外，总结下来有以下几点：

\begin{itemize}
\tightlist
\item
  HashMap允许一个null的key，value也可以为null；但是HashTable不允许null作为key或者value
\item
  HashMap是JDK1.2才引入的
\item
  HashTable是基于Dictionary接口实现的
\end{itemize}

HashTable（以及ConcurrentHashMap）都是不允许null值作为Key和Vaule的，主要的原因是因为要支持并发，假设调用\texttt{get(key)}得到了null，你是不能确认是key不存在，还是说存在但是值为null。在非并发场景下可以通过\texttt{contains(key)}来判断是否真的存在，但是在并发场景下，很可能会被其他线程修改。在JDK注释中有这样的解释：

\begin{quote}
The main reason that nulls aren't allowed in ConcurrentMaps (ConcurrentHashMaps, ConcurrentSkipListMaps) is that ambiguities that may be just barely tolerable in non-concurrent maps can't be accommodated. The main one is that if map.get(key) returns null, you can't detect whether the key explicitly maps to null vs the key isn't mapped. In a non-concurrent map, you can check this via map.contains(key), but in a concurrent one, the map might have changed between calls.
\end{quote}

\section{ConcurrentHashMap}\label{concurrenthashmap}

在Java1.7和1.8中ConcurrentHashMap实现差别较大，在1.7中采用分段锁的方式实现，将Map分为许多个Segment（Segment继承自ReentrantLock）,操作的时候，只会去占用某一个Segment，而其他的Segment不会受到影响。

而在1.8中直接使用CAS+ synchronized来实现。其在内存中的结构与HashMap几乎相同了。

\subsection{get操作}\label{getux64cdux4f5c}

\begin{Shaded}
\begin{Highlighting}[]
\KeywordTok{public}\NormalTok{ V }\FunctionTok{get}\OperatorTok{(}\BuiltInTok{Object}\NormalTok{ key}\OperatorTok{)} \OperatorTok{\{}
    \BuiltInTok{Node}\OperatorTok{\textless{}}\NormalTok{K}\OperatorTok{,}\NormalTok{V}\OperatorTok{\textgreater{}[]}\NormalTok{ tab}\OperatorTok{;} \BuiltInTok{Node}\OperatorTok{\textless{}}\NormalTok{K}\OperatorTok{,}\NormalTok{V}\OperatorTok{\textgreater{}}\NormalTok{ e}\OperatorTok{,}\NormalTok{ p}\OperatorTok{;} \DataTypeTok{int}\NormalTok{ n}\OperatorTok{,}\NormalTok{ eh}\OperatorTok{;}\NormalTok{ K ek}\OperatorTok{;}

    \CommentTok{// 得到最终的hash值（将高位混合到低位去避免哈希冲突）}
    \DataTypeTok{int}\NormalTok{ h }\OperatorTok{=} \FunctionTok{spread}\OperatorTok{(}\NormalTok{key}\OperatorTok{.}\FunctionTok{hashCode}\OperatorTok{());}
    \ControlFlowTok{if} \OperatorTok{((}\NormalTok{tab }\OperatorTok{=}\NormalTok{ table}\OperatorTok{)} \OperatorTok{!=} \KeywordTok{null} \OperatorTok{\&\&} \OperatorTok{(}\NormalTok{n }\OperatorTok{=}\NormalTok{ tab}\OperatorTok{.}\FunctionTok{length}\OperatorTok{)} \OperatorTok{\textgreater{}} \DecValTok{0} \OperatorTok{\&\&}
        \CommentTok{// (n{-}1) \& h 计算出所在的index值（与HashMap相同）}
        \OperatorTok{(}\NormalTok{e }\OperatorTok{=} \FunctionTok{tabAt}\OperatorTok{(}\NormalTok{tab}\OperatorTok{,} \OperatorTok{(}\NormalTok{n }\OperatorTok{{-}} \DecValTok{1}\OperatorTok{)} \OperatorTok{\&}\NormalTok{ h}\OperatorTok{))} \OperatorTok{!=} \KeywordTok{null}\OperatorTok{)} \OperatorTok{\{} 
        \CommentTok{// 如果哈希值相同，则直接定位到节点，再判断是否equal即可}
        \ControlFlowTok{if} \OperatorTok{((}\NormalTok{eh }\OperatorTok{=}\NormalTok{ e}\OperatorTok{.}\FunctionTok{hash}\OperatorTok{)} \OperatorTok{==}\NormalTok{ h}\OperatorTok{)} \OperatorTok{\{}                
            \CommentTok{// 这里比较key是否equal，单纯只凭hashCode是不够的。}
            \CommentTok{// 首先比较内存地址是否一致；然后再调用equals方法，是一种优化手段。}
            \ControlFlowTok{if} \OperatorTok{((}\NormalTok{ek }\OperatorTok{=}\NormalTok{ e}\OperatorTok{.}\FunctionTok{key}\OperatorTok{)} \OperatorTok{==}\NormalTok{ key }\OperatorTok{||} \OperatorTok{(}\NormalTok{ek }\OperatorTok{!=} \KeywordTok{null} \OperatorTok{\&\&}\NormalTok{ key}\OperatorTok{.}\FunctionTok{equals}\OperatorTok{(}\NormalTok{ek}\OperatorTok{)))}
                \ControlFlowTok{return}\NormalTok{ e}\OperatorTok{.}\FunctionTok{val}\OperatorTok{;}
        \OperatorTok{\}}
        \CommentTok{/*}
\CommentTok{            Hash值的首位被用作标记位，为负数的hash值是特殊的节点（也就是红黑树化了）}
\CommentTok{        */}
        \CommentTok{// 如果根据哈希值没有匹配到，那证明可能有哈希冲突，为负数是红黑树则在树中查找}
        \ControlFlowTok{else} \ControlFlowTok{if} \OperatorTok{(}\NormalTok{eh }\OperatorTok{\textless{}} \DecValTok{0}\OperatorTok{)}                         
            \ControlFlowTok{return} \OperatorTok{(}\NormalTok{p }\OperatorTok{=}\NormalTok{ e}\OperatorTok{.}\FunctionTok{find}\OperatorTok{(}\NormalTok{h}\OperatorTok{,}\NormalTok{ key}\OperatorTok{))} \OperatorTok{!=} \KeywordTok{null} \OperatorTok{?}\NormalTok{ p}\OperatorTok{.}\FunctionTok{val} \OperatorTok{:} \KeywordTok{null}\OperatorTok{;}
        \CommentTok{// 否则是普通的链表，在链表中一直朝下找即可}
        \ControlFlowTok{while} \OperatorTok{((}\NormalTok{e }\OperatorTok{=}\NormalTok{ e}\OperatorTok{.}\FunctionTok{next}\OperatorTok{)} \OperatorTok{!=} \KeywordTok{null}\OperatorTok{)} \OperatorTok{\{}           
            \ControlFlowTok{if} \OperatorTok{(}\NormalTok{e}\OperatorTok{.}\FunctionTok{hash} \OperatorTok{==}\NormalTok{ h }\OperatorTok{\&\&}
                \OperatorTok{((}\NormalTok{ek }\OperatorTok{=}\NormalTok{ e}\OperatorTok{.}\FunctionTok{key}\OperatorTok{)} \OperatorTok{==}\NormalTok{ key }\OperatorTok{||} \OperatorTok{(}\NormalTok{ek }\OperatorTok{!=} \KeywordTok{null} \OperatorTok{\&\&}\NormalTok{ key}\OperatorTok{.}\FunctionTok{equals}\OperatorTok{(}\NormalTok{ek}\OperatorTok{))))}
                \ControlFlowTok{return}\NormalTok{ e}\OperatorTok{.}\FunctionTok{val}\OperatorTok{;}
        \OperatorTok{\}}
    \OperatorTok{\}}
    \ControlFlowTok{return} \KeywordTok{null}\OperatorTok{;}
\OperatorTok{\}}
\end{Highlighting}
\end{Shaded}

可见get操作没有加任何的锁，而是通过将\texttt{transient\ volatile\ Node\textless{}K,V\textgreater{}{[}{]}\ table;}将table设置为volatile来保证可见性的。

\begin{Shaded}
\begin{Highlighting}[]
\KeywordTok{transient} \KeywordTok{volatile} \BuiltInTok{Node}\OperatorTok{\textless{}}\NormalTok{K}\OperatorTok{,}\NormalTok{V}\OperatorTok{\textgreater{}[]}\NormalTok{ table}\OperatorTok{;}

\DataTypeTok{static} \KeywordTok{class} \BuiltInTok{Node}\OperatorTok{\textless{}}\NormalTok{K}\OperatorTok{,}\NormalTok{V}\OperatorTok{\textgreater{}} \KeywordTok{implements} \BuiltInTok{Map}\OperatorTok{.}\FunctionTok{Entry}\OperatorTok{\textless{}}\NormalTok{K}\OperatorTok{,}\NormalTok{V}\OperatorTok{\textgreater{}} \OperatorTok{\{}
    \DataTypeTok{final} \DataTypeTok{int}\NormalTok{ hash}\OperatorTok{;}
    \DataTypeTok{final}\NormalTok{ K key}\OperatorTok{;}
    \KeywordTok{volatile}\NormalTok{ V val}\OperatorTok{;}
    \KeywordTok{volatile} \BuiltInTok{Node}\OperatorTok{\textless{}}\NormalTok{K}\OperatorTok{,}\NormalTok{V}\OperatorTok{\textgreater{}}\NormalTok{ next}\OperatorTok{;}
    \CommentTok{//...}
\OperatorTok{\}}
\end{Highlighting}
\end{Shaded}

\subsection{put操作}\label{putux64cdux4f5c}

\begin{Shaded}
\begin{Highlighting}[]
\DataTypeTok{final}\NormalTok{ V }\FunctionTok{putVal}\OperatorTok{(}\NormalTok{K key}\OperatorTok{,}\NormalTok{ V value}\OperatorTok{,} \DataTypeTok{boolean}\NormalTok{ onlyIfAbsent}\OperatorTok{)} \OperatorTok{\{}
    \ControlFlowTok{if} \OperatorTok{(}\NormalTok{key }\OperatorTok{==} \KeywordTok{null} \OperatorTok{||}\NormalTok{ value }\OperatorTok{==} \KeywordTok{null}\OperatorTok{)} \ControlFlowTok{throw} \KeywordTok{new} \BuiltInTok{NullPointerException}\OperatorTok{();}
    \DataTypeTok{int}\NormalTok{ hash }\OperatorTok{=} \FunctionTok{spread}\OperatorTok{(}\NormalTok{key}\OperatorTok{.}\FunctionTok{hashCode}\OperatorTok{());}
    \DataTypeTok{int}\NormalTok{ binCount }\OperatorTok{=} \DecValTok{0}\OperatorTok{;}
    \ControlFlowTok{for} \OperatorTok{(}\BuiltInTok{Node}\OperatorTok{\textless{}}\NormalTok{K}\OperatorTok{,}\NormalTok{V}\OperatorTok{\textgreater{}[]}\NormalTok{ tab }\OperatorTok{=}\NormalTok{ table}\OperatorTok{;;)} \OperatorTok{\{}
        \BuiltInTok{Node}\OperatorTok{\textless{}}\NormalTok{K}\OperatorTok{,}\NormalTok{V}\OperatorTok{\textgreater{}}\NormalTok{ f}\OperatorTok{;} \DataTypeTok{int}\NormalTok{ n}\OperatorTok{,}\NormalTok{ i}\OperatorTok{,}\NormalTok{ fh}\OperatorTok{;}
        \CommentTok{// 因为是懒加载，第一次插入的时候可能需要初始化}
        \ControlFlowTok{if} \OperatorTok{(}\NormalTok{tab }\OperatorTok{==} \KeywordTok{null} \OperatorTok{||} \OperatorTok{(}\NormalTok{n }\OperatorTok{=}\NormalTok{ tab}\OperatorTok{.}\FunctionTok{length}\OperatorTok{)} \OperatorTok{==} \DecValTok{0}\OperatorTok{)}
\NormalTok{            tab }\OperatorTok{=} \FunctionTok{initTable}\OperatorTok{();}
        \CommentTok{// 没有找到（节点之前不存在）}
        \ControlFlowTok{else} \ControlFlowTok{if} \OperatorTok{((}\NormalTok{f }\OperatorTok{=} \FunctionTok{tabAt}\OperatorTok{(}\NormalTok{tab}\OperatorTok{,}\NormalTok{ i }\OperatorTok{=} \OperatorTok{(}\NormalTok{n }\OperatorTok{{-}} \DecValTok{1}\OperatorTok{)} \OperatorTok{\&}\NormalTok{ hash}\OperatorTok{))} \OperatorTok{==} \KeywordTok{null}\OperatorTok{)} \OperatorTok{\{}
            \CommentTok{// 尝试CAS插入节点到空桶中，如果失败，则会重新走上面的流程进来}
            \ControlFlowTok{if} \OperatorTok{(}\FunctionTok{casTabAt}\OperatorTok{(}\NormalTok{tab}\OperatorTok{,}\NormalTok{ i}\OperatorTok{,} \KeywordTok{null}\OperatorTok{,}
                         \KeywordTok{new} \BuiltInTok{Node}\OperatorTok{\textless{}}\NormalTok{K}\OperatorTok{,}\NormalTok{V}\OperatorTok{\textgreater{}(}\NormalTok{hash}\OperatorTok{,}\NormalTok{ key}\OperatorTok{,}\NormalTok{ value}\OperatorTok{,} \KeywordTok{null}\OperatorTok{)))}
                \ControlFlowTok{break}\OperatorTok{;}                   \CommentTok{// no lock when adding to empty bin}
        \OperatorTok{\}}
        \CommentTok{// 如果（后面的流程）插入到了红黑树中，会导致首节点改变，所以这个地方需要帮忙更改过来}
        \ControlFlowTok{else} \ControlFlowTok{if} \OperatorTok{((}\NormalTok{fh }\OperatorTok{=}\NormalTok{ f}\OperatorTok{.}\FunctionTok{hash}\OperatorTok{)} \OperatorTok{==}\NormalTok{ MOVED}\OperatorTok{)}
\NormalTok{            tab }\OperatorTok{=} \FunctionTok{helpTransfer}\OperatorTok{(}\NormalTok{tab}\OperatorTok{,}\NormalTok{ f}\OperatorTok{);}
        \ControlFlowTok{else} \OperatorTok{\{}
\NormalTok{            V oldVal }\OperatorTok{=} \KeywordTok{null}\OperatorTok{;}
            \CommentTok{// f为当前定位到的桶中的第一个节点，将其同步进行后续操作}
            \KeywordTok{synchronized} \OperatorTok{(}\NormalTok{f}\OperatorTok{)} \OperatorTok{\{}
                \CommentTok{// 看看同步之前当前节点是否已经被更改了；如果是则需要重新开始轮回}
                \ControlFlowTok{if} \OperatorTok{(}\FunctionTok{tabAt}\OperatorTok{(}\NormalTok{tab}\OperatorTok{,}\NormalTok{ i}\OperatorTok{)} \OperatorTok{==}\NormalTok{ f}\OperatorTok{)} \OperatorTok{\{}
                    \CommentTok{// 普通链表}
                    \ControlFlowTok{if} \OperatorTok{(}\NormalTok{fh }\OperatorTok{\textgreater{}=} \DecValTok{0}\OperatorTok{)} \OperatorTok{\{}
\NormalTok{                        binCount }\OperatorTok{=} \DecValTok{1}\OperatorTok{;}
                        \ControlFlowTok{for} \OperatorTok{(}\BuiltInTok{Node}\OperatorTok{\textless{}}\NormalTok{K}\OperatorTok{,}\NormalTok{V}\OperatorTok{\textgreater{}}\NormalTok{ e }\OperatorTok{=}\NormalTok{ f}\OperatorTok{;;} \OperatorTok{++}\NormalTok{binCount}\OperatorTok{)} \OperatorTok{\{}
\NormalTok{                            K ek}\OperatorTok{;}
                            \CommentTok{// 如果已经存在值}
                            \ControlFlowTok{if} \OperatorTok{(}\NormalTok{e}\OperatorTok{.}\FunctionTok{hash} \OperatorTok{==}\NormalTok{ hash }\OperatorTok{\&\&}
                                \OperatorTok{((}\NormalTok{ek }\OperatorTok{=}\NormalTok{ e}\OperatorTok{.}\FunctionTok{key}\OperatorTok{)} \OperatorTok{==}\NormalTok{ key }\OperatorTok{||}
                                 \OperatorTok{(}\NormalTok{ek }\OperatorTok{!=} \KeywordTok{null} \OperatorTok{\&\&}\NormalTok{ key}\OperatorTok{.}\FunctionTok{equals}\OperatorTok{(}\NormalTok{ek}\OperatorTok{))))} \OperatorTok{\{}
\NormalTok{                                oldVal }\OperatorTok{=}\NormalTok{ e}\OperatorTok{.}\FunctionTok{val}\OperatorTok{;}
                                \ControlFlowTok{if} \OperatorTok{(!}\NormalTok{onlyIfAbsent}\OperatorTok{)}
\NormalTok{                                    e}\OperatorTok{.}\FunctionTok{val} \OperatorTok{=}\NormalTok{ value}\OperatorTok{;}
                                \ControlFlowTok{break}\OperatorTok{;}
                            \OperatorTok{\}}
                            \CommentTok{// 不存在则新增一个节点，插入到链表尾部}
                            \BuiltInTok{Node}\OperatorTok{\textless{}}\NormalTok{K}\OperatorTok{,}\NormalTok{V}\OperatorTok{\textgreater{}}\NormalTok{ pred }\OperatorTok{=}\NormalTok{ e}\OperatorTok{;}
                            \ControlFlowTok{if} \OperatorTok{((}\NormalTok{e }\OperatorTok{=}\NormalTok{ e}\OperatorTok{.}\FunctionTok{next}\OperatorTok{)} \OperatorTok{==} \KeywordTok{null}\OperatorTok{)} \OperatorTok{\{}
\NormalTok{                                pred}\OperatorTok{.}\FunctionTok{next} \OperatorTok{=} \KeywordTok{new} \BuiltInTok{Node}\OperatorTok{\textless{}}\NormalTok{K}\OperatorTok{,}\NormalTok{V}\OperatorTok{\textgreater{}(}\NormalTok{hash}\OperatorTok{,}\NormalTok{ key}\OperatorTok{,}
\NormalTok{                                                          value}\OperatorTok{,} \KeywordTok{null}\OperatorTok{);}
                                \ControlFlowTok{break}\OperatorTok{;}
                            \OperatorTok{\}}
                        \OperatorTok{\}}
                    \OperatorTok{\}}
                    \CommentTok{// 按红黑树处理}
                    \ControlFlowTok{else} \ControlFlowTok{if} \OperatorTok{(}\NormalTok{f }\KeywordTok{instanceof}\NormalTok{ TreeBin}\OperatorTok{)} \OperatorTok{\{}
                        \BuiltInTok{Node}\OperatorTok{\textless{}}\NormalTok{K}\OperatorTok{,}\NormalTok{V}\OperatorTok{\textgreater{}}\NormalTok{ p}\OperatorTok{;}
\NormalTok{                        binCount }\OperatorTok{=} \DecValTok{2}\OperatorTok{;}
                        \ControlFlowTok{if} \OperatorTok{((}\NormalTok{p }\OperatorTok{=} \OperatorTok{((}\NormalTok{TreeBin}\OperatorTok{\textless{}}\NormalTok{K}\OperatorTok{,}\NormalTok{V}\OperatorTok{\textgreater{})}\NormalTok{f}\OperatorTok{).}\FunctionTok{putTreeVal}\OperatorTok{(}\NormalTok{hash}\OperatorTok{,}\NormalTok{ key}\OperatorTok{,}
\NormalTok{                                                       value}\OperatorTok{))} \OperatorTok{!=} \KeywordTok{null}\OperatorTok{)} \OperatorTok{\{}
\NormalTok{                            oldVal }\OperatorTok{=}\NormalTok{ p}\OperatorTok{.}\FunctionTok{val}\OperatorTok{;}
                            \ControlFlowTok{if} \OperatorTok{(!}\NormalTok{onlyIfAbsent}\OperatorTok{)}
\NormalTok{                                p}\OperatorTok{.}\FunctionTok{val} \OperatorTok{=}\NormalTok{ value}\OperatorTok{;}
                        \OperatorTok{\}}
                    \OperatorTok{\}}
                \OperatorTok{\}}
            \OperatorTok{\}}
            \ControlFlowTok{if} \OperatorTok{(}\NormalTok{binCount }\OperatorTok{!=} \DecValTok{0}\OperatorTok{)} \OperatorTok{\{}
                \ControlFlowTok{if} \OperatorTok{(}\NormalTok{binCount }\OperatorTok{\textgreater{}=}\NormalTok{ TREEIFY\_THRESHOLD}\OperatorTok{)}
                    \FunctionTok{treeifyBin}\OperatorTok{(}\NormalTok{tab}\OperatorTok{,}\NormalTok{ i}\OperatorTok{);}
                \ControlFlowTok{if} \OperatorTok{(}\NormalTok{oldVal }\OperatorTok{!=} \KeywordTok{null}\OperatorTok{)}
                    \ControlFlowTok{return}\NormalTok{ oldVal}\OperatorTok{;}
                \ControlFlowTok{break}\OperatorTok{;}
            \OperatorTok{\}}
        \OperatorTok{\}}
    \OperatorTok{\}}
    \FunctionTok{addCount}\OperatorTok{(}\DecValTok{1L}\OperatorTok{,}\NormalTok{ binCount}\OperatorTok{);}
    \ControlFlowTok{return} \KeywordTok{null}\OperatorTok{;}
\OperatorTok{\}}
\end{Highlighting}
\end{Shaded}

\begin{itemize}
\tightlist
\item
  \href{https://www.javatpoint.com/difference-between-hashmap-and-hashtable}{Difference between HashMap and Hashtable}
\end{itemize}


\end{document}
